\documentclass{scrartcl}
\usepackage{graphicx}  % required for inserting images

% for cyrillic and Bulgarian language
\usepackage[utf8]{inputenc}
\usepackage[T2A]{fontenc}
\usepackage[bulgarian]{babel}

% math packages
\usepackage{amsmath}
\usepackage{float}
\usepackage{amsfonts}
\usepackage{amssymb}
\usepackage{amsthm}
\usepackage{cancel}

\usepackage{ragged2e}  % adds FlushLeft

\usepackage{listings}
\usepackage{xcolor}

\definecolor{codegreen}{rgb}{0,0.6,0}
\definecolor{codegray}{rgb}{0.5,0.5,0.5}
\definecolor{codepurple}{rgb}{0.58,0,0.82}
\definecolor{backcolour}{rgb}{0.95,0.95,0.92}

\lstdefinestyle{mystyle}{
    language=Octave,
    backgroundcolor=\color{backcolour},   
    commentstyle=\color{codegreen},
    keywordstyle=\color{magenta},
    numberstyle=\tiny\color{codegray},
    stringstyle=\color{codepurple},
    basicstyle=\ttfamily\footnotesize,
    breakatwhitespace=false,         
    breaklines=true,
    captionpos=b,
    keepspaces=false,                 
    numbersep=5pt,
    showspaces=false,
    showstringspaces=false,
    showtabs=false,
    tabsize=2,
    columns=fullflexible,
}

\lstset{style=mystyle}

\usepackage[colorlinks=true, linkcolor=blue, urlcolor=cyan]{hyperref}

\newtheorem{theorem}{Теорема}
\newtheorem{definition}{Дефиниция}
\newtheorem{sledstvie}{Следствие}

\title{Линейна алгебра - собствени вектори и собствени стойности}
\author{Ивайло Андреев}
\date{14 май 2026}

\begin{document}
\maketitle

\tableofcontents

\newpage

\begin{FlushLeft}

\section{Определение на характеристичен полином на матрица}

\begin{definition}
(Характеристичен полином и характеристични корени на матрица)

Нека $A \in M_n(F)$ е квадратна матрица.

Характеристичен полином на матрицата $A$ e

$$
P_A(x) = \det(A - x E) =
\begin{vmatrix}
    a_{11} - x & a_{12} & \dots & a_{1n}\\
    a_{21} & a_{22} - x & \dots & a_{2n}\\
    \vdots & \vdots & \ddots & \vdots\\
    a_{n1} & a_{n2} & \dots & a_{nn} - x
\end{vmatrix}
$$

Характеристични корени на матрицата $A$ са корените на $P_A(x) = 0$.
\end{definition}

От разлагането на детерминантата чрез пермутации можем да запишем характеристичния полином по следния начин:

$$P_A(x) = (-1)^n x^n + (-1)^{n-1}(\text{tr }A) x^{n-1}+ \dots + \det A$$

където $\text{tr }A = \displaystyle\sum_{i=1}^{n}a_{ii}$ е следата на матрицата $A$, която е сумата на елементите по главния диагонал.

Ако характеристичните корени на матрицата $A$ са $\lambda_1, \dots, \lambda_n$ (не непременно различни), то можем да запишем характеристичния полином по следния начин:

$$P_A(x) = (-1)^n \prod_{i=1}^n (x - \lambda_i)$$

\section{Подобни матрици и доказателство, че подобните матрици имат едни и същи характеристични полиноми}

\begin{definition}
(Подобна матрица)

Нека $A, B \in M_n(F)$.

Матриците $A$ и $B$ са подобни (и пишем $A \sim B$), ако съществува обратима матрица $T \in M_n(F)$, такава че

$$B = T^{-1}AT$$
\end{definition}

\begin{theorem}\label{th:sim_matrix}
Подобните матрици имат един и същи характеристичен полином.
\end{theorem}

\textbf{Доказателство:}

Нека $A, B \in M_n(F)$ и нека $A \sim B$. Тогава съществува обратима матрица $T$, такава че $B = T^{-1}AT$

Ще покажем, че $P_B(x) = P_A(x)$.

\begin{align*}
P_B(x)
&= \det(B - x E)\\
&= \det(T^{-1}AT - x E)\\
&= \det(T^{-1}AT - x T^{-1} ET)\\
&= \det(T^{-1}(A - x E)T)\\
&= \det(T^{-1})\det(A - x E)\det(T)\\
&= \det(T)^{-1}\det(A - x E)\det(T)\\
&= \det(A - x E)\\
&= P_A(x)\\
\end{align*}

\begin{sledstvie}
Нека $V$ е линейно пространство над числово поле $F$ и нека $\dim V = n < \infty$.
Нека $\varphi \in \text{Hom }{V}$ е линеен оператор.
Нека $\{e_i\}_{i=1}^n$ и $\{h_i\}_{i=1}^n$ са базиси на $V$.
Нека $A$ е матрицата на $\varphi$ в базиса $\{e_i\}_{i=1}^n$ и нека $B$ е матрицата на $\varphi$ в базиса $\{h_i\}_{i=1}^n$.

Знаем, че $A \sim B$, защото са матрици на $\varphi$ в различни базиси. От \ref{th:sim_matrix} следва, че

$$P_A(x) = P_B(x)$$

Казано по-свободно - матриците на един линеен оператор $\varphi \in \text{Hom }V$ в различни базиси на $V$ притежават едни и същи характеристични полиноми.
\end{sledstvie}

Това следствие ни позволява да дефинираме характеристичен полином на линеен оператор.

\section{Определение на характеристичен полином на линеен оператор}

\begin{definition}
(Характеристичен полином и характеристични корени на линеен оператор)

Нека $V$ е линейно пространство над числово поле $F$ и нека $\dim V = n < \infty$.
Нека $\varphi \in \text{Hom }{V}$ е линеен оператор.
Нека $\{e_i\}_{i=1}^n$ е базис на $V$ и нека $A$ е матрицата на $\varphi$ в базиса $\{e_i\}_{i=1}^n$.

Тогава характеристичен полином на линейния оператор $\varphi$ е характеристичния полином на матрицата на $\varphi$ в базиса $\{e_i\}_{i=1}^n$, тоест

$$P_{\varphi}(x) = P_A(x)$$

Характеристични корени на линейния оператор $\varphi$ са корените на $P_{\varphi}(x) = 0$.
\end{definition}

\section{Определение на
собствен вектор и собствена стойност на линеен оператор}

\begin{definition}
(Собствен вектор и собствена стойност)

Нека $V$ е линейно пространство над числово поле $F$ и нека $\dim V = n < \infty$.
Нека $\varphi \in \text{Hom }{V}$ е линеен оператор.
Нека $v\in V$.

Наричаме вектора $v$ собствен вектор, съответстващ на собствената стойност $\lambda \in F$, ако

\begin{enumerate}
    \item $v \ne \mathbf{0}$
    \item $\varphi(v) = \lambda v$
\end{enumerate}
\end{definition}

\section{Необходимо и достатъчно
условие характеристичен корен на линеен оператор да е собствена стойност в
крайномерно линейно пространство}

\begin{theorem}
Нека $V$ е линейно пространство над числово поле $F$ и нека $\dim V = n < \infty$.
Нека $\varphi \in \text{Hom }{V}$ е линеен оператор.

Тогава характеристични корени на линейния оператор $\varphi$, които са в полето $F$, и само те, са всички собствени стойности на линейния оператор $\varphi$.
\end{theorem}

\textbf{Доказателство:}

Нека $\{e_i\}_{i=1}^n$ е фиксиран базис на $V$.

Нека $A$ е матрицата на $\varphi$ в базиса $\{e_i\}_{i=1}^n$.

\textbf{Права посока}

Нека $\lambda \in F$ е собствена стойност, съответстваща на ненулевия собствен вектор $v \neq \mathbf{0}$. Ще покажем, че $\lambda$ е корен на характеристичния полином на $\varphi$.

По дефиниция

$$\varphi(v) = \lambda v$$

Нека координатите на $v$ в базиса $\{e_i\}_{i=1}^n$ са $\alpha = \begin{pmatrix}
    \alpha_1\\
    \vdots\\
    \alpha_n
\end{pmatrix}$

$$v = \sum_{i=1}^n \alpha_i e_i$$

Знаем, че $v$ е собствен вектор, в частност ненулев, и съответно координатите му не са нулев вектор. Тоест

$$\alpha \ne \mathbf{0}$$

Нека координатите на $\varphi(v)$ в базиса $\{e_i\}_{i=1}^n$ са $\beta = \begin{pmatrix}
    \beta_1\\
    \vdots\\
    \beta_n
\end{pmatrix}$

$$\varphi(v) = \sum_{i=1}^n \beta_i e_i$$

Знаем, че $v$ е собствен вектор и е изпълнено равенството

$$\varphi(v) = \lambda v$$

$$\sum_{i=1}^n \beta_i e_i = \lambda \sum_{i=1}^n \alpha_i e_i$$

$$\sum_{i=1}^n \beta_i e_i = \sum_{i=1}^n (\lambda\alpha_i) e_i$$

От единствеността на представянето по базис следва, че

$$\beta = \lambda \alpha$$

От друга страна от връзката между координатните стълбове и матрицата на линейния оператор следва, че

$$\beta = A \alpha$$

Приравняваме и получаваме

$$\beta = A \alpha = \lambda \alpha$$

$$A \alpha - \lambda \alpha = 0$$

$$A \alpha - \lambda E \alpha = 0$$

$$(A  - \lambda E) \alpha = 0$$

Получаваме хомогенна система линейни уравнения. Вече споменахме, че $\alpha \ne \mathbf{0}$ и съответно системата притежава нетривиални решения. Тогава детерминантата на матрицата $A - \lambda E$ е нула.

$$\det{(A - \lambda E)} = 0$$

$$P_A(\lambda) = 0$$

Така получаваме, че $\lambda$, която по начало казахме, че е собствена стойност на линейния оператор, също е и корен на характеристичния полином на линейния оператор.

\textbf{Обратна посока}

Нека $\lambda \in F$ е корен на характеристичния полином на линейния оператор $\varphi$. Ще покажем, че $\lambda$ е собствена стойност на $\varphi$.

По дефиниция

$$\det(A - \lambda E) = 0$$

Разглеждаме хомогенната система линейни уравнения спрямо неизвестния вектор-стълб $x$

$$(A - \lambda E)x = \mathbf{0}$$

Тъй като детерминантата на матрицата на системата е нула, системата притежава ненулево решение $\alpha \neq \mathbf{0} = \begin{pmatrix}
    \alpha_1\\
    \vdots\\
    \alpha_n
\end{pmatrix}$, такова че

$$A \alpha = \lambda \alpha$$

Нека $v = \sum_{i=1}^n \alpha_i e_i$ е векторът от $V$, чийто координати в базиса $\{e_i\}_{i=1}^n$ са точно елементите на вектора $\alpha$. Знаем, че $\alpha \neq \mathbf{0}$, то и $v \neq \mathbf{0}$.

От матричното представяне на линейния оператор знаем, че ако координатите на $\varphi(v)$ в базиса $\{e_i\}_{i=1}^n$ са $\beta = \begin{pmatrix}
    \beta_1\\
    \vdots\\
    \beta_n
\end{pmatrix}$, то е изпълнено равенството 

$$\beta = A \alpha$$

След приравняване

$$\beta = A \alpha = \lambda \alpha$$

Това означава, че векторът $\varphi(v)$ има координати $\lambda \alpha$, т.е.

$$\varphi(v) = \sum_{i=1}^n (\lambda \alpha_i) e_i = \lambda \sum_{i=1}^n \alpha_i e_i = \lambda v$$

Тъй като $v \neq \mathbf{0}$ и $\varphi(v) = \lambda v$, то по дефиниция $\lambda$ е собствена стойност на линейния оператор $\varphi$.

С това теоремата е доказана в двете посоки.

\section{Доказателство, че собствените вектори на линеен
оператор, съответстващи на различни собствени стойности са линейно независими}

\begin{theorem}\label{theorem_linear_independency}
Нека $V$ е линейно пространство над числово поле $F$ и нека $\dim V = k < \infty$.
Нека $\varphi \in \text{Hom }{V}$ е линеен оператор.

Тогава собствените вектори на $\varphi$, съответстващи на различни собствени стойности, са линейно независими.
\end{theorem}

\textbf{Доказателство:}

Нека $\lambda_1, \dots, \lambda_k$ са различни собствени стойности, съответстващи на собствените вектори $v_1, \dots, v_k$.

Индукция по $k$.

\textbf{База}: $k= 1$

По определение собствения вектор е ненулев и знаем, че система от един ненулев вектор е линейно независима.

\textbf{Индукционно предположение}: $k = n$

Нека $\lambda_1, \dots, \lambda_n$ са различни собствени стойности, съответстващи на линейно независимите собствени вектори $v_1, \dots, v_n$. Допускаме, че твърдението е вярно за всяка система от $n$ вектора.

\textbf{Стъпка}: $k = n + 1$

Разглеждаме равенството

\begin{equation}
    \beta_1 v_1 + \dots + \beta_n v_n + \beta_{n+1} v_{n+1} = \mathbf{0} \label{before_phi}
\end{equation}

Ще покажем, че всички коефициенти $\beta_i = 0 \quad i \in \{1 \dots n+1\}$ от линейната комбинация са нули, откъдето ще следва, че векторите са линейно независими.

$$\varphi(\beta_1 v_1 + \dots + \beta_n v_n + \beta_{n+1} v_{n+1}) = \varphi(\mathbf{0})$$

$$\varphi(\beta_1 v_1) + \dots + \varphi(\beta_n v_n) + \varphi(\beta_{n+1} v_{n+1}) = \varphi(\mathbf{0})$$

\begin{equation}
    \lambda_1\beta_1 v_1 + \dots + \lambda_n\beta_n v_n + \lambda_{n+1}\beta_{n+1} v_{n+1} = \mathbf{0} \label{after_phi}
\end{equation}

Имаме уравненията \eqref{before_phi} и \eqref{after_phi}. Целим се да елиминираме терма $v_{n+1}$, за да можем да стигнем до уравнение, където да можем да използваме индукционното предположение. Ще постигнем това като умножим равенството \eqref{before_phi} с $\lambda_{n+1}$ и го извадим от равенството \eqref{after_phi}.

$$(\lambda_1 - \lambda_{n+1})\beta_1 v_1 + \dots + (\lambda_n - \lambda_{n+1})\beta_n v_n + (\lambda_{n+1} - \lambda_{n+1})\beta_{n+1} v_{n+1} = \mathbf{0}$$

\begin{equation}
    (\lambda_1 - \lambda_{n+1})\beta_1 v_1 + \dots + (\lambda_n - \lambda_{n+1})\beta_n v_n = \mathbf{0} \label{after_subtraction}
\end{equation}

От индукционното предположение знаем, че $v_1, \dots, v_n$ са линейно независими. Това означава, че коефициентите в линейната комбинация, даваща нула, също са нули. В \eqref{after_subtraction} коефициентите $\lambda_i - \lambda_{n+1} \ne 0 \quad i \in \{1 \dots n\}$ са ненулеви, защото собствените стойности са различни по условие. Тогава остава $\beta_i = 0 \quad i \in \{1 \dots n\}$.

Заместваме в \eqref{before_phi} с нулевите коефициенти

$$\beta_{n+1} v_{n+1} = \mathbf{0}$$

Векторът $v_{n+1}$ е собствен и в частност ненулев, откъдето $\beta_{n+1} = 0$.

Показахме, че $\beta_i = 0 \quad i \in \{1 \dots n+1\}$, с което теоремата е доказана.

\section{Линеен оператор с прост спектър (с доказателство)}

\begin{definition}
(Линеен оператор с прост спектър)

Нека $V$ е линейно пространство над числово поле $F$ и нека $\dim V = n < \infty$.
Нека $\varphi \in \text{Hom }{V}$ е линеен оператор.

Тогава линейния оператор $\varphi$ има прост спектър, ако има $n$ на брой различни характеристични корени.
\end{definition}

\begin{theorem}
Нека $V$ е линейно пространство над числово поле $F$ и нека $\dim V = n < \infty$.
Нека $\varphi \in \text{Hom }{V}$ е линеен оператор с прост спектър.

Тогава съществува базис на $V$, в който матрицата на $\varphi$ е диагонална.
\end{theorem}

\textbf{Доказателство:}

Линейният оператор $\varphi$ е с прост спектър и съответно притежава $n$ на брой различни собствени стойности. Нека това са $\lambda_1$, $\dots$, $\lambda_n$. Нека на тези собствени стойности съответстват собствените вектори $v_1$, $\dots$, $v_n$. Според теорема \ref{theorem_linear_independency} $v_1$, $\dots$, $v_n$ са линейно независими и тъй като са $n$ на брой, образуват базис на $V$.

Знаем, че $\varphi(v_i) = \lambda_i v_i$, понеже векторите са собствени. От определението за матрица на линеен оператор следва, че матрицата на $\varphi$ в базиса на собствените вектори има вида

$$
D =
\begin{pmatrix}
    \lambda_1 & & 0\\
    & \ddots &\\
    0 & & \lambda_n
\end{pmatrix}
$$

Матрицата $D$ е диагонална, с което теоремата е доказана.

\section{Определение за ф-инвариантно
подпространство на линейно пространство}

\begin{definition}
(Инвариантно подпространство)

Нека $V$ е линейно пространство над числово поле $F$.
Нека $\varphi \in \text{Hom }{V}$ е линеен оператор.
Нека $U$ е подпространство на $V$.

Тогава подпространството $U$ e $\varphi$-инвариантно, ако за всеки вектор $u \in U$ е изпълнено, че $\varphi(u) \in U$.
\end{definition}

\section{Основни примери за ф-инвариантни
подпространства (с доказателство)}

\subsection{Пример 1}

Очевидно $U = V$ е $\varphi$-инвариантно подпространство.

$$\dots$$

Подпространството $U = \lbrace \mathbf{0} \rbrace$ е $\varphi$-инвариантно за всеки линеен оператор $\varphi$, защото линейните оператори имат свойството, че $\varphi(\mathbf{0}) = \mathbf{0}$.

\subsection{Пример 2}

Подпространството $\text{Ker}\varphi$ е $\varphi$-инвариантно подпространство.

Нека $u \in \text{Ker}\varphi$. Тогава $\varphi(u) = \mathbf{0}$

$$\varphi(\varphi(u)) = \varphi(\mathbf{0}) = \mathbf{0} \in \text{Ker}\varphi$$

$$\dots$$

Подпространството $\text{Im}\varphi$ е $\varphi$-инвариантно подпространство.

Нека $b \in \text{Im}\varphi$. Тогава съществува $a$, за който $\varphi(a) = b$.

$$\varphi(\varphi(a)) = \varphi(b) \in \text{Im}\varphi$$

\subsection{Пример 3}

Подпространството $U = \{v \in V \space | \space \varphi(v) = \lambda v\}$, където $\lambda$ е фиксирана собствена стойност, е $\varphi$-инвариантно подпространство.

Нека $u \in U$. Тогава $\varphi(u) = \lambda u$.

$$\varphi(\varphi(u)) = \varphi(\lambda u) = \lambda\varphi(u)$$

Нека означим $w = \varphi(u)$. Тогава

$$\varphi(w) = \lambda w$$

Следователно $w \in U$.

$$\dots$$

Подпространството $U = l(v)$, където $v$ е фиксиран собствен вектор, е $\varphi$-инвариантно подпространство.

Нека $u \in U$. Тогава съществува скалар $\alpha \in F$, такъв че $u = \alpha v$.

Знаем, че $v$ е собствен вектор и съответно $\varphi(v) = \lambda v$.

$$\varphi(u) = \varphi(\alpha v) = \alpha\varphi(v) = \alpha(\lambda v) = (\alpha\lambda)v \in U$$

\section{Теореми за съществуване на ф-инвариантни
подпространства на линейно пространство над полето на реалните и на комплексните числа (без доказателство)}

\begin{theorem}
Нека $V$ е линейно пространство над $\mathbb{R}$ и нека $\dim{V} = n < \infty$ и нека $\varphi \in \text{Hom }{V}$. Тогава $\varphi$ има едномерно или двумерно $\varphi$-инвариантно подпространство.
\end{theorem}

\begin{theorem}
Нека $V$ е линейно пространство над $\mathbb{C}$ и нека $\dim{V} = n < \infty$ и нека $\varphi \in \text{Hom }{V}$. Тогава за всяка собствена стойност $\lambda \in \mathbb{C}$ на $\varphi$ съществува едномерно $\varphi$-инвариантно подпространство, което е линейната обвивка на някой собствен вектор, съответстващ на $\lambda$.
\end{theorem}

\end{FlushLeft}

\end{document}